\documentclass[11pt,a4paper,oneside]{book}
\usepackage{mathscommun}
% Les présentations, packages et environnements pédagogiques proviennent de mathscommun.sty.
% Ces alias préservent uniquement les couleurs des figures du cours historique.
\usetikzlibrary{babel}
\colorlet{bleuNuit}{encre}
\colorlet{bleu}{accent}
\colorlet{bleuPale}{grisclair}
\colorlet{vert}{accent}
\colorlet{vertPale}{grisclair}
\colorlet{orangePale}{grisclair}
\colorlet{gris}{black!55}
\colorlet{grisPale}{grisclair}
\newcommand{\niveau}[1]{{\footnotesize\sffamily\color{gris} #1}}

\setniveau{Spécialité : trigonométrie}
\renewcommand{\ReferenceProgramme}{Dossier thématique : seconde, première, terminale et prolongements}
\renewcommand{\EditionDocument}{Édition du 31 octobre 2026}

\hypersetup{pdftitle={Spécialité trigonométrie - Cours et démonstrations}}
\begin{document}
\titrecollection{Cours et démonstrations}{Comprendre les fondements avant les formules}{De la longueur du cercle et de \ensuremath{\pi} aux radians, rotations, identités, équations, dérivées et applications. Les preuves et les résultats admis sont distingués.}
\frontmatter
\tableofcontents
\mainmatter
\subfile{chapitres/01_origine_pi.tex}
\subfile{chapitres/02_acquis_geometriques.tex}
\subfile{chapitres/03_radians.tex}
\subfile{chapitres/04_cercle_unite.tex}
\subfile{chapitres/05_valeurs_remarquables.tex}
\subfile{chapitres/06_symetries_periodes.tex}
\subfile{chapitres/07_addition_rotations.tex}
\subfile{chapitres/08_identites.tex}
\subfile{chapitres/09_tangente.tex}
\subfile{chapitres/10_courbes.tex}
\subfile{chapitres/11_equations.tex}
\subfile{chapitres/12_triangles.tex}
\subfile{chapitres/13_vecteurs.tex}
\subfile{chapitres/14_analyse.tex}
\subfile{chapitres/15_reciproques.tex}
\subfile{chapitres/16_oscillations.tex}
\subfile{chapitres/17_complexes.tex}
\subfile{chapitres/18_methodes.tex}
\subfile{chapitres/19_synthese.tex}
\end{document}
